\documentclass[11pt, a4paper]{article}

% --- Pacotes Fundamentais ---
\usepackage[utf8]{inputenc}
\usepackage[T1]{fontenc}
\usepackage[brazil]{babel}
\usepackage{microtype}
\usepackage{geometry}
\usepackage{titlesec}
\usepackage{fancyhdr}
\usepackage{xcolor}
\usepackage{booktabs}
\usepackage{tabularx}
\usepackage{float}
\usepackage{parskip}
\usepackage{lmodern}
\usepackage{enumitem}
\usepackage{hyperref}
\usepackage{lastpage}

% --- Definição de Cores ---
\definecolor{primaryColor}{RGB}{0, 51, 102}
\definecolor{secondaryText}{RGB}{80, 80, 80}
\definecolor{accentColor}{RGB}{178, 34, 34}

% --- Configuração de Hyperlinks ---
\hypersetup{
    colorlinks=false,
    pdftitle={Contrato de Prestação de Serviços — Desenvolvimento de Software},
    pdfauthor={Info Tecnologia e Ecossistema em Saúde Ltda}
}

% --- Layout de Página ---
\geometry{left=2.5cm, right=2.5cm, top=3.8cm, bottom=2.2cm, headsep=1.2cm, headheight=50pt}

% --- Cabeçalho e Rodapé ---
\pagestyle{fancy}
\fancyhf{}

\fancyhead[C]{%
    \begin{tabularx}{\textwidth}{@{} X r @{}}
        \textcolor{primaryColor}{\textbf{\footnotesize CONTRATO DE PRESTAÇÃO DE SERVIÇOS — DESENVOLVIMENTO DE SOFTWARE}} & \scriptsize \textit{Info Dental} \\
        \textcolor{secondaryText}{\scriptsize CNPJ 48.818.166/0001-46 \,|\, Info Tecnologia e Ecossistema em Saúde Ltda} & {\scriptsize \textbf{Abril 2026}}
    \end{tabularx}%
}

\fancyfoot[L]{\scriptsize \textcolor{secondaryText}{Info Dental \,|\, Documento confidencial}}
\fancyfoot[R]{\normalsize \textbf{\thepage/\pageref{LastPage}}}

\renewcommand{\headrulewidth}{1.5pt}
\renewcommand{\headrule}{%
    \vspace{0.1cm}
    \hbox to\headwidth{\color{primaryColor}\leaders\hrule height \headrulewidth\hfill}%
}

\renewcommand{\footrulewidth}{1.5pt}
\renewcommand{\footrule}{%
    \color{primaryColor}%
    \hbox to\headwidth{\leaders\hrule height \footrulewidth\hfill}%
    \vspace{0.2cm}
}

% --- Formatação de Títulos ---
\titleformat{\subsection}
{\color{primaryColor}\normalfont\normalsize\bfseries}
{}{0em}{}
\titlespacing*{\subsection}{0pt}{10pt}{5pt}

% --- Hifenização e justificação ---
\tolerance=1000
\emergencystretch=1.5em
\hyphenpenalty=50
\exhyphenpenalty=50

% --- Compactar listas ---
\setlist{nosep, leftmargin=1.2em}

% --- INÍCIO DO DOCUMENTO ---
\begin{document}

% =============================================
% TÍTULO
% =============================================
\begin{center}
    \vspace*{0.3cm}
    {\color{primaryColor}\rule{\textwidth}{2pt}}
    \vspace{0.4cm}

    {\LARGE\bfseries\color{primaryColor} CONTRATO DE PRESTAÇÃO DE SERVIÇOS\\[4pt] PROFISSIONAIS ESPECIALIZADOS EM\\[4pt] DESENVOLVIMENTO DE SOFTWARE}

    \vspace{0.4cm}
    {\color{primaryColor}\rule{\textwidth}{2pt}}
\end{center}

\vspace{0.8cm}

% =============================================
% QUALIFICAÇÃO DAS PARTES
% =============================================

\textcolor{primaryColor}{\textbf{CONTRATANTE:}}

\textbf{INFO TECNOLOGIA E ECOSSISTEMA EM SAÚDE LTDA}, pessoa jurídica de direito privado, inscrita no CNPJ sob o nº 48.818.166/0001-46, com sede na Rua 05, nº 691, Loja 01-B — Edif.\ The Prime Tamandaré, Setor Oeste, Goiânia/GO, CEP 74115-060, neste ato representada por seu administrador \textbf{KRONER MACHADO COSTA}, brasileiro, portador do CPF nº 029.808.481-31, residente e domiciliado na Rua 08, nº 200, Apartamento 602, Residencial Caiobá, Setor Oeste, Goiânia-GO, CEP 74115-100.

\vspace{0.4cm}

\textcolor{primaryColor}{\textbf{CONTRATADA:}}

\textbf{66.107.449 WALKER FREITAS DOS SANTOS}, Microempreendedor Individual, inscrito no CNPJ sob o nº \textbf{66.107.449/0001-61}, com sede na Alameda dos Buritis, nº 408, Quadra 83, Lote 02, Pavmto Superior, Setor Central, Goiânia/GO, CEP 74015-080, neste ato representado por \textbf{WALKER FREITAS DOS SANTOS}, brasileiro, solteiro, portador da carteira de identidade RG nº 7760636, expedida pela PC/PA em 18/07/2024, inscrito no CPF sob o nº 040.169.402-06, natural de Redenção/PA, residente e domiciliado no endereço supracitado.

\vspace{0.5cm}

As partes acima identificadas têm, entre si, justo e acertado o presente Contrato de Prestação de Serviços Profissionais Especializados em Desenvolvimento de Software, que se regerá pelas cláusulas e condições seguintes:

% --- CLÁUSULA PRIMEIRA ---
\subsection{CLÁUSULA PRIMEIRA — DO OBJETO DO CONTRATO}

1.1.\ O presente contrato tem por objeto a prestação de serviços especializados na área de \textbf{desenvolvimento de software} pela CONTRATADA em favor da CONTRATANTE, compreendendo, mas não se limitando a:

\begin{enumerate}[label=\alph*), leftmargin=1.8em]
    \item Desenvolvimento, manutenção e evolução de sistemas e aplicações de software, incluindo front-end, back-end, APIs, integrações e microsserviços;
    \item Implementação de funcionalidades, correção de bugs, otimização de performance e refatoração de código;
    \item Desenvolvimento e manutenção de aplicações mobile (Android e/ou iOS), quando aplicável;
    \item Criação e manutenção de testes automatizados (unitários, de integração e end-to-end);
    \item Elaboração de documentação técnica, incluindo especificações de APIs, diagramas de arquitetura, fluxos de dados e manuais de implantação;
    \item Configuração e manutenção de ambientes de desenvolvimento, homologação e produção, incluindo infraestrutura em nuvem, containers (Docker) e pipelines de CI/CD;
    \item Participação em revisões de código (code reviews), definição de padrões de desenvolvimento e boas práticas;
    \item Outras atividades correlatas à função de desenvolvedor de software, conforme demandas específicas da CONTRATANTE, previamente alinhadas entre as partes.
\end{enumerate}

1.2.\ Os serviços serão prestados pela CONTRATADA com total autonomia técnica e sem qualquer vínculo empregatício com a CONTRATANTE, não havendo subordinação jurídica, pessoalidade ou habitualidade nos moldes da Consolidação das Leis do Trabalho (CLT).

% --- CLÁUSULA SEGUNDA ---
\subsection{CLÁUSULA SEGUNDA — DA EXECUÇÃO DOS SERVIÇOS}

2.1.\ A CONTRATADA organizará suas atividades e horários para o cumprimento do objeto deste contrato, comprometendo-se a atender às demandas e prazos acordados com a CONTRATANTE. Reuniões de alinhamento, planejamento de sprints e acompanhamento de projetos poderão ser agendadas conforme necessidade mútua.

2.2.\ Os serviços poderão ser prestados remotamente, a partir das instalações da CONTRATADA, ou, eventualmente, nas dependências da CONTRATANTE, conforme previamente acordado entre as partes.

2.3.\ A CONTRATADA deverá utilizar os repositórios de código, ferramentas de gerenciamento de projetos e canais de comunicação indicados pela CONTRATANTE, comprometendo-se a seguir os fluxos de trabalho, padrões de versionamento (branching strategy) e processos de deploy definidos pela CONTRATANTE.

2.4.\ Todo código produzido pela CONTRATADA deverá ser commitado nos repositórios oficiais da CONTRATANTE, sendo vedado o armazenamento de código-fonte em repositórios pessoais ou de terceiros.

% --- CLÁUSULA TERCEIRA ---
\subsection{CLÁUSULA TERCEIRA — DAS OBRIGAÇÕES DA CONTRATANTE}

3.1.\ Fornecer à CONTRATADA todas as informações, requisitos funcionais e técnicos, documentação de referência, acesso a repositórios de código, ambientes de desenvolvimento/homologação/produção e demais recursos necessários à correta e eficiente execução dos serviços.

3.2.\ Efetuar o pagamento dos valores devidos à CONTRATADA, nos termos e condições estabelecidos na Cláusula Oitava.

3.3.\ Comunicar à CONTRATADA, com antecedência razoável, quaisquer alterações de escopo, requisitos, prioridades ou cronogramas que possam impactar a execução dos serviços.

3.4.\ Designar um ponto de contato principal (Product Owner, Tech Lead ou equivalente) para comunicação, priorização de demandas e aprovações técnicas junto à CONTRATADA.

3.5.\ Fornecer licenças de software e ferramentas necessárias à execução dos serviços, quando estas forem de uso obrigatório definido pela CONTRATANTE e não estiverem disponíveis em versão gratuita ou open source.

% --- CLÁUSULA QUARTA ---
\subsection{CLÁUSULA QUARTA — DAS OBRIGAÇÕES DA CONTRATADA}

4.1.\ Prestar os serviços descritos na Cláusula Primeira com zelo, diligência, qualidade técnica e em conformidade com as melhores práticas de desenvolvimento de software, incluindo padrões de clean code, segurança da informação e performance.

4.2.\ Emitir a correspondente Nota Fiscal de Serviços Eletrônica (NFS-e) referente aos valores recebidos, em conformidade com a legislação municipal e federal vigente. A não apresentação da Nota Fiscal poderá ensejar a retenção do pagamento até sua regularização.

4.3.\ Assumir integral responsabilidade por todos os encargos tributários (federais, estaduais e municipais), previdenciários e trabalhistas (caso possua empregados próprios) decorrentes da execução deste contrato e de sua atividade empresarial, isentando a CONTRATANTE de qualquer responsabilidade nesse sentido.

4.4.\ Manter seus dados cadastrais atualizados junto à CONTRATANTE.

4.5.\ Utilizar os equipamentos, ferramentas e softwares necessários para a prestação dos serviços, salvo aqueles cuja responsabilidade de fornecimento seja da CONTRATANTE nos termos do item 3.5.

4.6.\ Reportar à CONTRATANTE o andamento dos projetos, incluindo progresso de tarefas, impedimentos técnicos e estimativas de entrega, por meio dos canais e ferramentas de gerenciamento definidos pela CONTRATANTE.

4.7.\ Produzir código-fonte legível, documentado e aderente aos padrões de codificação estabelecidos pela CONTRATANTE, realizando commits atômicos e descritivos.

4.8.\ Entregar à CONTRATANTE, ao término de cada ciclo de desenvolvimento ou mediante solicitação, toda a documentação técnica atualizada referente aos módulos e funcionalidades desenvolvidos.

4.9.\ Não introduzir, deliberada ou negligentemente, vulnerabilidades de segurança, backdoors, código malicioso ou quaisquer elementos que possam comprometer a integridade, disponibilidade ou confidencialidade dos sistemas da CONTRATANTE.

% --- CLÁUSULA QUINTA ---
\subsection{CLÁUSULA QUINTA — DA CONFIDENCIALIDADE E SIGILO}

5.1.\ A CONTRATADA obriga-se a manter sigilo absoluto sobre todas as informações a que tiver acesso em razão da execução deste contrato, incluindo, de forma exemplificativa e não exaustiva:

\begin{enumerate}[label=\alph*), leftmargin=1.8em]
    \item Código-fonte, algoritmos, estruturas de dados, arquitetura de sistemas e modelagem de bancos de dados dos softwares da CONTRATANTE;
    \item Credenciais de acesso a servidores, APIs, bancos de dados, repositórios, plataformas de cloud computing e quaisquer outros sistemas;
    \item Regras de negócio, fluxos operacionais e lógica de processamento implementados nos softwares;
    \item Dados de clientes, pacientes, usuários e quaisquer informações pessoais ou sensíveis armazenados ou processados pelos sistemas;
    \item Informações comerciais, estratégicas, financeiras, de pricing, roadmap de produto e planos de desenvolvimento futuro;
    \item Segredos industriais, know-how tecnológico, metodologias proprietárias e processos internos de desenvolvimento;
    \item Documentação técnica, diagramas, especificações, atas de reuniões e quaisquer registros relacionados aos projetos.
\end{enumerate}

5.2.\ A obrigação de sigilo \textbf{perdurará por prazo indeterminado}, mesmo após a extinção ou rescisão deste contrato, independentemente do motivo.

5.3.\ A CONTRATADA não poderá, em hipótese alguma, copiar, reproduzir, compartilhar, divulgar, transferir ou utilizar, para fins diversos do objeto contratual, quaisquer das informações confidenciais mencionadas no item 5.1, sob pena de responder civil e criminalmente pelas perdas e danos causados, sem prejuízo das penalidades previstas na Cláusula Décima Quinta.

5.4.\ A CONTRATADA deverá garantir que seus sócios, prepostos, subcontratados ou colaboradores que porventura venham a ter acesso às informações confidenciais da CONTRATANTE também observem o dever de sigilo aqui estabelecido, respondendo solidariamente por eventuais violações.

5.5.\ As obrigações de confidencialidade não se aplicam a informações que: (a) sejam ou se tornem públicas sem culpa da CONTRATADA; (b) tenham sido legitimamente conhecidas pela CONTRATADA antes da celebração deste contrato, mediante comprovação documental; (c) sejam determinadas por ordem judicial ou administrativa a serem reveladas, caso em que a CONTRATADA deverá notificar imediatamente a CONTRATANTE antes da divulgação.

5.6.\ Ao término ou rescisão deste contrato, por qualquer motivo, a CONTRATADA deverá, no prazo máximo de \textbf{5 (cinco) dias úteis}, devolver à CONTRATANTE ou destruir de forma segura e irrecuperável todas as cópias de informações confidenciais, código-fonte, documentação, credenciais e quaisquer materiais em seu poder, independentemente do meio ou suporte em que se encontrem (digital, impresso, em nuvem, dispositivos pessoais, etc.), emitindo declaração por escrito confirmando o cumprimento desta obrigação.

% --- CLÁUSULA SEXTA ---
\subsection{CLÁUSULA SEXTA — DA PROPRIEDADE INTELECTUAL}

6.1.\ Nos termos do art.\ 4º, \textit{caput}, da Lei nº 9.609/98 (Lei do Software), \textbf{todos os direitos patrimoniais} sobre os programas de computador, sistemas, módulos, funcionalidades, bibliotecas proprietárias, scripts, automações, integrações, documentação técnica, especificações, diagramas de arquitetura, modelagens de banco de dados e quaisquer outras obras intelectuais desenvolvidas pela CONTRATADA no âmbito deste contrato (``Desenvolvimentos'') \textbf{serão de titularidade exclusiva e integral da CONTRATANTE}.

6.2.\ A CONTRATADA \textbf{cede e transfere} à CONTRATANTE, em caráter \textbf{irrevogável e irretratável}, todos os direitos patrimoniais de autor sobre os Desenvolvimentos, podendo a CONTRATANTE utilizá-los, modificá-los, licenciá-los, sublicenciá-los, distribuí-los e comercializá-los livremente, por prazo indeterminado, em qualquer meio ou território, sem qualquer ônus adicional.

6.3.\ A cessão de que trata esta cláusula abrange, de forma exemplificativa: código-fonte, código compilado, interfaces, layouts, bases de dados (estrutura e conteúdo gerado), documentação, manuais, materiais de treinamento e quaisquer derivações ou evoluções dos Desenvolvimentos.

6.4.\ A compensação pelo trabalho intelectual da CONTRATADA \textbf{limita-se exclusivamente à remuneração} prevista na Cláusula Oitava deste contrato, nos termos do art.\ 4º, \textit{caput}, parte final, da Lei nº 9.609/98, não fazendo jus a qualquer remuneração adicional a título de royalties, participação nos lucros ou qualquer outra forma de compensação.

6.5.\ A CONTRATADA declara e garante que os Desenvolvimentos serão originais e não infringirão direitos de propriedade intelectual de terceiros, incluindo patentes, direitos autorais, segredos industriais ou marcas, responsabilizando-se integralmente por quaisquer reclamações, ações ou demandas nesse sentido, inclusive isentando a CONTRATANTE de eventuais condenações.

6.6.\ Caso a CONTRATADA utilize, nos Desenvolvimentos, bibliotecas ou componentes de terceiros licenciados sob código aberto (open source), deverá informar previamente à CONTRATANTE, identificando a licença aplicável, para que a CONTRATANTE avalie a compatibilidade com seus produtos e estratégia de negócio. É vedado o uso de componentes open source com licenças copyleft (GPL, AGPL e similares) sem aprovação expressa e por escrito da CONTRATANTE.

6.7.\ A CONTRATADA não poderá registrar ou reivindicar, em seu nome ou de terceiros, quaisquer direitos de propriedade intelectual sobre os Desenvolvimentos.

% --- CLÁUSULA SÉTIMA ---
\subsection{CLÁUSULA SÉTIMA — DA NÃO CONCORRÊNCIA}

7.1.\ Durante a vigência deste contrato e pelo prazo de \textbf{2 (dois) anos} contados da sua extinção ou rescisão, por qualquer motivo, a CONTRATADA obriga-se a \textbf{não exercer}, direta ou indiretamente, por conta própria ou por intermédio de terceiros, atividades que configurem concorrência com a CONTRATANTE, incluindo, de forma exemplificativa:

\begin{enumerate}[label=\alph*), leftmargin=1.8em]
    \item Desenvolver, criar, comercializar, licenciar, distribuir ou prestar serviços para empresas ou projetos que envolvam \textbf{software de gestão para clínicas odontológicas ou de saúde} que concorra diretamente com os produtos da CONTRATANTE;
    \item Prestar serviços de consultoria, assessoria técnica ou desenvolvimento de software para empresas concorrentes diretas da CONTRATANTE no segmento de tecnologia para saúde/odontologia;
    \item Participar, como sócio, acionista, administrador, consultor, empregado ou prestador de serviços, de empresas que atuem no mesmo segmento de mercado da CONTRATANTE e que com ela concorram diretamente.
\end{enumerate}

7.2.\ A restrição de não concorrência prevista nesta cláusula aplica-se a todo o \textbf{território nacional brasileiro e aos Estados Unidos da América}, abrangendo os mercados em que a CONTRATANTE atua ou venha a atuar durante a vigência deste contrato.

7.3.\ As restrições desta cláusula \textbf{não impedem} a CONTRATADA de prestar serviços de desenvolvimento de software em segmentos de mercado distintos daquele em que atua a CONTRATANTE, desde que não envolva o uso ou divulgação de informações confidenciais obtidas em razão deste contrato.

7.4.\ O descumprimento da obrigação de não concorrência sujeitará a CONTRATADA ao pagamento de multa não compensatória no valor de \textbf{12 (doze) vezes o valor mensal} dos serviços vigente à época da rescisão, sem prejuízo do dever de indenizar por perdas e danos efetivamente comprovados.

% --- CLÁUSULA OITAVA ---
\subsection{CLÁUSULA OITAVA — DO VALOR E DA FORMA DE PAGAMENTO}

8.1.\ Pelos serviços prestados, a CONTRATANTE pagará à CONTRATADA o valor mensal de \textbf{R\$ 3.500,00 (três mil e quinhentos reais)}.

8.2.\ O pagamento será efetuado até o \textbf{dia 10 (dez) do mês subsequente} à prestação dos serviços, mediante depósito/transferência bancária ou PIX para a conta de titularidade da CONTRATADA, cujos dados são:

\begin{itemize}[label=\textcolor{primaryColor}{$\bullet$}, leftmargin=1.8em]
    \item \textbf{Banco:} 077 — Inter
    \item \textbf{Agência:} 0001
    \item \textbf{Conta Corrente:} 52663919-9
    \item \textbf{CNPJ (para PIX):} 66.107.449/0001-61
\end{itemize}

8.3.\ O primeiro pagamento será realizado \textit{pro rata die}, caso o início da prestação dos serviços ocorra após o primeiro dia do mês.

% --- CLÁUSULA NONA ---
\subsection{CLÁUSULA NONA — DA VIGÊNCIA}

9.1.\ O presente contrato entra em vigor na data de sua assinatura e terá validade pelo prazo de \textbf{12 (doze) meses}, podendo ser prorrogado automaticamente por iguais períodos, caso não haja manifestação contrária de qualquer das partes com antecedência mínima de 30 (trinta) dias do término do período de vigência.

% --- CLÁUSULA DÉCIMA ---
\subsection{CLÁUSULA DÉCIMA — DA RESCISÃO}

10.1.\ Este contrato poderá ser rescindido por qualquer das partes, a qualquer tempo e imotivadamente, mediante comunicação por escrito à outra parte, com antecedência mínima de \textbf{30 (trinta) dias}.

10.1.1.\ Na hipótese de rescisão imotivada com aviso prévio inferior a 30 (trinta) dias, a parte que der causa à rescisão sem o cumprimento integral do aviso pagará à outra parte multa compensatória no valor correspondente aos dias faltantes do aviso prévio, calculado com base no valor mensal dos serviços.

10.2.\ O contrato será rescindido de pleno direito, independentemente de qualquer notificação ou interpelação judicial ou extrajudicial, nas seguintes hipóteses:

\begin{enumerate}[label=\alph*), leftmargin=1.8em]
    \item Descumprimento de quaisquer cláusulas ou condições aqui estabelecidas por qualquer das partes, desde que a parte infratora, notificada por escrito, não sanar a falha no prazo de 10 (dez) dias;
    \item Pedido de recuperação judicial, extrajudicial, falência ou insolvência civil de qualquer das partes;
    \item Violação das cláusulas de confidencialidade (Cláusula Quinta), propriedade intelectual (Cláusula Sexta) ou não concorrência (Cláusula Sétima) pela CONTRATADA, hipótese em que a rescisão operará imediatamente, sem necessidade de prazo para cura.
\end{enumerate}

10.3.\ Em caso de rescisão, a CONTRATANTE efetuará o pagamento dos serviços efetivamente prestados pela CONTRATADA até a data da efetiva rescisão, de forma proporcional.

% --- CLÁUSULA DÉCIMA PRIMEIRA ---
\subsection{CLÁUSULA DÉCIMA PRIMEIRA — DA TRANSIÇÃO E ENTREGA TÉCNICA}

11.1.\ Em caso de rescisão ou término deste contrato, por qualquer motivo, a CONTRATADA obriga-se a realizar, durante o período de aviso prévio, a \textbf{transição técnica completa} dos projetos e atividades em andamento, incluindo:

\begin{enumerate}[label=\alph*), leftmargin=1.8em]
    \item Entrega de todo código-fonte atualizado nos repositórios oficiais da CONTRATANTE, com todas as branches e trabalhos em progresso devidamente commitados;
    \item Atualização de toda a documentação técnica dos módulos e funcionalidades desenvolvidos ou em desenvolvimento;
    \item Realização de sessões de transferência de conhecimento (knowledge transfer) com a equipe técnica da CONTRATANTE ou com o profissional designado para assumir as atividades;
    \item Entrega de todas as credenciais, acessos e configurações de ambientes sob responsabilidade da CONTRATADA;
    \item Devolução ou destruição de todos os materiais, dados e informações confidenciais, conforme Cláusula Quinta.
\end{enumerate}

11.2.\ A obrigação de transição técnica aplica-se mesmo em caso de rescisão por justa causa ou imediata, devendo a CONTRATADA colaborar para a continuidade dos serviços pelo prazo mínimo necessário à transferência, não inferior a \textbf{5 (cinco) dias úteis}, independentemente do aviso prévio.

% --- CLÁUSULA DÉCIMA SEGUNDA ---
\subsection{CLÁUSULA DÉCIMA SEGUNDA — DA NÃO SOLICITAÇÃO (NON-SOLICITATION)}

12.1.\ Durante a vigência deste contrato e por um período de \textbf{12 (doze) meses} após sua extinção ou rescisão, a CONTRATADA obriga-se a \textbf{não aliciar, recrutar, contratar ou induzir} ao desligamento quaisquer empregados, prestadores de serviço, sócios, colaboradores ou parceiros comerciais da CONTRATANTE.

12.2.\ A vedação de que trata o item 12.1 aplica-se também à solicitação, direta ou indireta, de clientes da CONTRATANTE para oferecimento de serviços ou produtos concorrentes.

% --- CLÁUSULA DÉCIMA TERCEIRA ---
\subsection{CLÁUSULA DÉCIMA TERCEIRA — DA AUSÊNCIA DE VÍNCULO EMPREGATÍCIO}

13.1.\ Reafirmam as partes que o presente contrato é de natureza estritamente civil de prestação de serviços, não se estabelecendo qualquer vínculo empregatício, societário ou de representação comercial entre a CONTRATANTE e a CONTRATADA, seus sócios, administradores, empregados ou prepostos.

13.2.\ A CONTRATADA atuará com total autonomia na execução dos serviços, não estando sujeita à subordinação hierárquica, controle de jornada ou fiscalização direta da CONTRATANTE quanto à forma de execução das atividades, desde que observados os padrões de qualidade e os prazos acordados.

% --- CLÁUSULA DÉCIMA QUARTA ---
\subsection{CLÁUSULA DÉCIMA QUARTA — DA LEI GERAL DE PROTEÇÃO DE DADOS (LGPD)}

14.1.\ A CONTRATADA, no exercício do objeto deste contrato, terá acesso e realizará tratamento de dados pessoais cuja responsabilidade primária é da CONTRATANTE, atuando na qualidade de \textbf{Operadora}, nos termos da Lei nº 13.709/2018 (LGPD).

14.2.\ A CONTRATADA se compromete a tratar os dados pessoais \textbf{estritamente} para as finalidades objeto deste contrato, seguindo as instruções lícitas da CONTRATANTE, e a adotar medidas de segurança técnicas e administrativas aptas a proteger os dados pessoais de acessos não autorizados e de situações acidentais ou ilícitas de destruição, perda, alteração, comunicação ou qualquer forma de tratamento inadequado ou ilícito.

14.3.\ A CONTRATADA deverá implementar as seguintes medidas mínimas de segurança no tratamento de dados pessoais acessados em razão deste contrato:

\begin{enumerate}[label=\alph*), leftmargin=1.8em]
    \item Utilizar conexões seguras (VPN, SSH, HTTPS) para acesso a ambientes que contenham dados pessoais;
    \item Não armazenar dados pessoais de clientes ou usuários do software em ambientes locais ou pessoais;
    \item Utilizar criptografia para transmissão e armazenamento de dados sensíveis;
    \item Reportar imediatamente à CONTRATANTE qualquer suspeita de acesso não autorizado ou incidente de segurança.
\end{enumerate}

14.4.\ A CONTRATADA manterá registro das operações de tratamento de dados pessoais que realizar e comunicará à CONTRATANTE, \textbf{em até 24 (vinte e quatro) horas}, a ocorrência de qualquer incidente de segurança que possa acarretar risco ou dano relevante aos titulares dos dados.

14.5.\ Ao término deste contrato, ou mediante solicitação da CONTRATANTE, a CONTRATADA deverá eliminar de forma segura ou devolver todos os dados pessoais a que teve acesso, salvo se a manutenção for exigida por obrigação legal.

% --- CLÁUSULA DÉCIMA QUINTA ---
\subsection{CLÁUSULA DÉCIMA QUINTA — DAS PENALIDADES}

15.1.\ O descumprimento das obrigações de confidencialidade (Cláusula Quinta) sujeitará a CONTRATADA ao pagamento de multa no valor de \textbf{R\$ 50.000,00 (cinquenta mil reais)} por evento de violação, sem prejuízo da obrigação de indenizar as perdas e danos efetivamente causados.

15.2.\ O descumprimento das obrigações de propriedade intelectual (Cláusula Sexta), incluindo a reivindicação de direitos sobre os Desenvolvimentos ou a utilização não autorizada de código-fonte da CONTRATANTE, sujeitará a CONTRATADA ao pagamento de multa no valor de \textbf{R\$ 100.000,00 (cem mil reais)}, sem prejuízo das medidas judiciais cabíveis, inclusive de natureza criminal, nos termos da Lei nº 9.609/98.

15.3.\ O descumprimento da obrigação de não concorrência (Cláusula Sétima) sujeitará a CONTRATADA à multa prevista no item 7.4, cumulável com indenização por perdas e danos.

15.4.\ As multas previstas nesta cláusula são cumuláveis entre si e com outras penalidades previstas neste contrato, e não excluem o direito da CONTRATANTE de buscar medidas judiciais de tutela de urgência (antecipação de tutela, medida cautelar) para cessar a violação.

% --- CLÁUSULA DÉCIMA SEXTA ---
\subsection{CLÁUSULA DÉCIMA SEXTA — DAS DISPOSIÇÕES GERAIS}

16.1.\ Eventuais alterações neste contrato somente serão válidas se realizadas por meio de Termo Aditivo escrito e assinado por ambas as partes.

16.2.\ A tolerância de uma das partes com relação ao descumprimento de qualquer cláusula ou condição deste contrato pela outra parte não implicará novação, renúncia ou alteração do pactuado, nem prejudicará o exercício do direito da parte tolerante de exigir o cumprimento da obrigação em momento posterior.

16.3.\ Este contrato obriga as partes e seus sucessores a qualquer título.

16.4.\ Nenhuma das partes poderá ceder ou transferir os direitos e obrigações decorrentes deste contrato a terceiros sem o prévio e expresso consentimento por escrito da outra parte.

16.5.\ A invalidade ou ineficácia de qualquer cláusula ou condição deste contrato não prejudicará a validade e eficácia das demais, que permanecerão em pleno vigor.

\pagebreak
% --- CLÁUSULA DÉCIMA SÉTIMA ---
\subsection{CLÁUSULA DÉCIMA SÉTIMA — DO FORO}

17.1.\ Fica eleito o foro da Comarca de Goiânia, Estado de Goiás, para dirimir quaisquer litígios ou dúvidas oriundas do presente contrato, com renúncia expressa a qualquer outro, por mais privilegiado que seja.

\vspace{1cm}

E, por estarem assim justas e contratadas, as partes assinam o presente instrumento em 2 (duas) vias de igual teor e forma, na presença das duas testemunhas abaixo.

\vspace{0.4cm}

\begin{center}
    Goiânia/GO, \_\_\_\_\_\_ de \_\_\_\_\_\_\_\_\_\_\_\_\_\_\_\_\_ de 2026.
\end{center}

\vspace{5.5cm}

\begin{center}
\begin{tabular}{p{0.45\textwidth} p{0.45\textwidth}}
    \centering\rule{7cm}{0.5pt} & \centering\rule{7cm}{0.5pt}\tabularnewline
    \centering\parbox[t]{7cm}{\centering\textbf{INFO TECNOLOGIA E ECOSSISTEMA\\[-2pt] EM SAÚDE LTDA}\\[2pt] \small Por: Kroner Machado Costa} &
    \centering\parbox[t]{7cm}{\centering\textbf{66.107.449 WALKER FREITAS\\[-2pt] DOS SANTOS}\\[2pt] \small Por: Walker Freitas dos Santos}\tabularnewline
\end{tabular}
\end{center}

\vspace{1.2cm}

\textbf{Testemunhas:}

\vspace{1cm}

\begin{tabular}{p{0.45\textwidth} p{0.45\textwidth}}
    1.\ Nome: \rule{5.5cm}{0.5pt} & 2.\ Nome: \rule{5.5cm}{0.5pt} \\[0.6cm]
    \hspace{0.35cm} CPF: \rule{5.5cm}{0.5pt} & \hspace{0.35cm} CPF: \rule{5.5cm}{0.5pt} \\
\end{tabular}

\end{document}
